% Source: 02 - Tue Sep 29 2026.pdf, page 19. Content remains in source order.
\sourcepage{02}{19}
\subsection{Class NP and the relationship between P and NP}

The complexity class NP is the set of all languages that are verified in polynomial time.
\[
\classNP=\{L\subseteq\bits^*:\ L\text{ verified in polynomial time}\}.
\]
NP can be equivalently defined with respect to decision problems : 
\[
\classNP=\{\Pi_C:\bits^*\to\{0,1\}:\ \Pi_C\text{ is verified in polynomial time}\}.
\]

We proved that $\mathrm{HAMILTON}\in\classNP$.

Since solving in polynomial time implies verifying in polynomial time (proof later), we have that 
\[
\text{if }L\in\classP\ \Rightarrow\ L\in\classNP.
\]
That is, $\classP\subseteq\classNP$.

But we do not know if $\classP=\classNP$ or $\classP\subset\classNP$.